\section{Density expansions of source gates}
\label{sec:peps_source_gate_density}

Fix the ordered pair slots of Definition~\ref{def:peps_source_slots},
with halfspaces $U_e,V_e$, and write $\mathcal S$ for their joint
register space. All tensor products below use this order and its
canonical associativity identifications. The slot set and all finite
index sets may be empty. These are the source-coordinate identities
used in \cite[proof of Theorem~5.2]{OpenAI2026PolynomialPEPS},
\texttt{04-compression.tex}, lines 279--355.

% This section proves exact preparation, matrix, and density identities.
% It does not assert that the actual Schmidt decomposition has yet been
% substituted into the finite-family identities. Gaussian moment estimates,
% simultaneous source replacement, and the circuit trace-norm error bound
% remain separate from these statements. No tag here certifies all of
% Theorem 5.2 of the source manuscript.

\begin{definition}[The joint source vector]
    \label{def:peps_joint_source_vector}
    \lean{TNLean.PEPS.PairEffect.SourceInventory.slotVector}
    \lean{TNLean.PEPS.PairEffect.SourceInventory.slotLayout_cons}
    \leanok
    \uses{def:peps_source_slots}
    The joint source vector is the multilinear map
    $\Phi:\prod_e(U_e\otimes V_e)\to\mathcal S$ defined recursively
    by $\Phi_{\varnothing}=1\in\C$ and
    $\Phi_{e::E}(\eta)=\eta_e\otimes\Phi_E(\eta|_E)$, with the
    first pair separated into its two registers. Thus $\Phi(\eta)$
    is the ordered tensor product of the pair vectors, including the
    terminal scalar factor.
\end{definition}

\begin{theorem}[Preparation and finite sums of sources]
    \label{thm:peps_source_preparation_sums}
    \lean{TNLean.PEPS.PairEffect.SourceInventory.slotVector_eq_vector}
    \lean{TNLean.PEPS.PairEffect.SourceInventory.eval_prepareSlots_eq_appendIso_symm}
    \lean{TNLean.PEPS.PairEffect.SourceInventory.eval_prepareSlots_sum_smul}
    \leanok
    \uses{def:peps_joint_source_vector}
    For every original memory $\mathcal H$, the actual source
    preparation satisfies
    $Q_{\eta,\mathcal H}x=\Phi(\eta)\otimes x$.
    In particular, let $J_e$ be finite, let $d_e:J_e\to\C$, and let
    $\theta_{e,j}\in U_e\otimes V_e$. Then
    \begin{align}
        Q_{(\sum_{j\in J_e}d_e(j)\theta_{e,j})_e,\mathcal H}
         & =\sum_{a\in\prod_e J_e}
        \left(\prod_e d_e(a_e)\right)
        Q_{(\theta_{e,a_e})_e,\mathcal H}.
        \label{eq:peps_source_preparation_sum}
    \end{align}
\end{theorem}

\begin{proof}\leanok
    \uses{def:peps_joint_source_vector,def:peps_source_inventory}
    Induction on the source list identifies $\Phi(\eta)$ with its
    actual joint vector: the empty list contributes $1$, and the
    induction step prepares $\eta_e$ in front of the remaining
    sources. Changing to the fixed register layout preserves this
    identity and its tensor product with $x$.
    Expand each argument of the multilinear map $\Phi$ as a finite
    sum. Its value on a scalar multiple in every argument is
    $\Phi((d_e\theta_e)_e)=(\prod_e d_e)\Phi((\theta_e)_e)$.
    Tensoring with $x$ gives (\ref{eq:peps_source_preparation_sum}).
\end{proof}

\begin{definition}[Matrices of prepared branches]
    \label{def:peps_prepared_matrix}
    \lean{TNLean.PEPS.PairEffect.Word.preparedMatrix}
    \leanok
    \uses{def:peps_source_slots,def:peps_party_layout}
    Let $T:\mathcal S\otimes\mathcal H_{\rm in}
        \to\mathcal H_{\rm out}$ be the operator of a source-only
    composition as in Definition~\ref{def:peps_party_layout}, without
    imposing normalization or contraction bounds.
    Choose finite orthonormal bases $f_j$ and $g_i$ of the input and
    output memories. Denote the matrix of the prepared branch by
    \begin{align}
        K_T(\eta)_{ij}
         & =\langle g_i,TQ_{\eta,\mathcal H_{\rm in}}f_j\rangle.
        \label{eq:peps_prepared_matrix}
    \end{align}
    The input and output bases are fixed throughout the density
    expansions below. No finite-dimensionality assumption on the
    private spaces is required for this definition.
\end{definition}

\begin{theorem}[Coordinates of source preparation]
    \label{thm:peps_prepared_matrix_coordinates}
    \lean{TNLean.PEPS.PairEffect.Word.preparedMatrix_sum_smul}
    \lean{TNLean.PEPS.PairEffect.Word.preparedMatrix_eq_sum_basis}
    \lean{TNLean.PEPS.PairEffect.SourceInventory.eval_prepareSlots_eq_sum_basis}
    \leanok
    \uses{def:peps_prepared_matrix}
    For finite families as in
    Theorem~\ref{thm:peps_source_preparation_sums},
    \begin{align}
        K_T\bigl((\sum_j d_e(j)\theta_{e,j})_e\bigr)
         & =\sum_{a\in\prod_e J_e}
        \left(\prod_e d_e(a_e)\right)K_T((\theta_{e,a_e})_e).
        \label{eq:peps_prepared_matrix_sum}
    \end{align}
    In particular, choose finite orthonormal bases $(u_{e,r})_r$ of
    $U_e$ and $(v_{e,s})_s$ of $V_e$, and put
    $\eta_e(r,s)=\langle u_{e,r}\otimes v_{e,s},\eta_e\rangle$.
    For $a_e=(r_e,s_e)$, set
    $\theta_{e,a_e}=u_{e,r_e}\otimes v_{e,s_e}$.
    Formula~(\ref{eq:peps_prepared_matrix_sum}) then holds with
    $d_e(a_e)=\eta_e(a_e)$ and left-hand side $K_T(\eta)$.
    The identical expansion holds for $Q_{\eta,\mathcal H}$.
\end{theorem}

\begin{proof}\leanok
    \uses{thm:peps_source_preparation_sums}
    Compose (\ref{eq:peps_source_preparation_sum}) with $T$ and take
    matrices in the prescribed bases. Both operations are linear.
    For the basis specialization, use
    $\eta_e=\sum_{r,s}\eta_e(r,s)u_{e,r}\otimes v_{e,s}$
    in each slot.
\end{proof}

\begin{definition}[Density coefficients of actual branches]
    \label{def:peps_prepared_density_coefficient}
    \lean{TNLean.PEPS.PairEffect.Word.preparedDensityCoefficient}
    \leanok
    \uses{def:peps_prepared_matrix}
    Choose vector families $u_{e,r}\in U_e$, $v_{e,s}\in V_e$ for
    a ket branch $T$, and independent families
    $\widetilde u_{e,r}\in U_e$, $\widetilde v_{e,s}\in V_e$
    for a bra branch $T'$. Their index sets may differ. Write
    $K_T(a)=K_T((u_{e,a_e^{(1)}}\otimes v_{e,a_e^{(2)}})_e)$
    and define $K_{T'}(b)$ using the second families.
    For any complex input matrix $\rho$, define
    \begin{align}
        C_{T,T'}(\rho;a,b) & =K_T(a)\rho K_{T'}(b)^\dagger.
        \label{eq:peps_prepared_density_coefficient}
    \end{align}
    The families need not span the private spaces or be orthogonal.
\end{definition}

\begin{theorem}[Density expansion for finite component families]
    \label{thm:peps_prepared_density_components}
    \lean{TNLean.PEPS.PairEffect.Word.preparedMatrix_density_sum_smul}
    \leanok
    \uses{def:peps_prepared_density_coefficient}
    Suppose the families in
    Definition~\ref{def:peps_prepared_density_coefficient} have finite
    index sets. Choose complex coefficients $d_e(r,s)$ and
    $\widetilde d_e(r,s)$, and put
    \begin{align}
        \eta_e  & =\sum_{r,s}d_e(r,s)u_{e,r}\otimes v_{e,s}, \\
        \zeta_e & =\sum_{r,s}\widetilde d_e(r,s)
        \widetilde u_{e,r}\otimes\widetilde v_{e,s}.
    \end{align}
    For every input matrix $\rho$,
    \begin{align}
        K_T(\eta)\rho K_{T'}(\zeta)^\dagger
         & =\sum_{a,b}\left(\prod_e
        d_e(a_e)\overline{\widetilde d_e(b_e)}\right)
        C_{T,T'}(\rho;a,b).
        \label{eq:peps_prepared_density_components}
    \end{align}
    This identity also applies to vector families contained in proper
    subspaces of the private halfspaces.
\end{theorem}

\begin{proof}\leanok
    \uses{thm:peps_prepared_matrix_coordinates}
    Expand the two prepared matrices using
    (\ref{eq:peps_prepared_matrix_sum}). The adjoint conjugates each
    scalar coefficient on the bra side. Distributing both finite sums
    yields (\ref{eq:peps_prepared_density_components}), since
    \begin{align*}
        \left(\prod_e d_e(a_e)\right)
        \overline{\prod_e\widetilde d_e(b_e)}
         & =\prod_e d_e(a_e)\overline{\widetilde d_e(b_e)}.
    \end{align*}
\end{proof}

\begin{theorem}[Coordinate density expansion of a weighted gate]
    \label{thm:peps_prepared_gate_density}
    \lean{TNLean.PEPS.PairEffect.Word.preparedMatrix_density_eq_sum_basis}
    \lean{TNLean.PEPS.PairEffect.Word.preparedMatrix_gate_density_eq_sum_basis}
    \leanok
    \uses{def:peps_prepared_density_coefficient}
    Let $\Xi$ be finite. For each $\xi\in\Xi$, choose a branch
    $T_\xi$, source vectors $\eta_{\xi,e}$, and finite orthonormal
    bases $(u_{\xi,e,r})_r$, $(v_{\xi,e,s})_s$ of the two halfspaces.
    The bases and their index sets may depend on $\xi$.
    Set $K_\xi=K_{T_\xi}(\eta_\xi)$ and
    $\kappa_\xi(a)=\prod_e\langle u_{\xi,e,a_e^{(1)}}\otimes
        v_{\xi,e,a_e^{(2)}},\eta_{\xi,e}\rangle$.
    For any input matrix $\rho$ and any two branches,
    \begin{align}
        K_\xi\rho K_\zeta^\dagger
         & =\sum_{a,b}\kappa_\xi(a)\overline{\kappa_\zeta(b)}
        C_{T_\xi,T_\zeta}(\rho;a,b).
        \label{eq:peps_prepared_density_basis}
    \end{align}
    If $G=\sum_\xi c_\xi K_\xi$, then
    \begin{align}
        G\rho G^\dagger
         & =\sum_{\xi,\zeta}c_\xi\overline{c_\zeta}
        \sum_{a,b}\kappa_\xi(a)\overline{\kappa_\zeta(b)}
        C_{T_\xi,T_\zeta}(\rho;a,b).
        \label{eq:peps_prepared_gate_density}
    \end{align}
    Neither identity requires positivity or normalization of $\rho$.
\end{theorem}

\begin{proof}\leanok
    \uses{thm:peps_prepared_matrix_coordinates}
    Expand each source in its branch's endpoint bases and distribute
    the ket and bra sums, conjugating only the bra coefficients. This
    gives (\ref{eq:peps_prepared_density_basis}). Expand
    $G\rho G^\dagger=\sum_{\xi,\zeta}
        c_\xi\overline{c_\zeta}K_\xi\rho K_\zeta^\dagger$ and substitute
    (\ref{eq:peps_prepared_density_basis}) to obtain
    (\ref{eq:peps_prepared_gate_density}).
\end{proof}

\begin{theorem}[Operator norm of a prepared branch]
    \label{thm:peps_prepared_matrix_contraction}
    \lean{TNLean.PEPS.PairEffect.Word.norm_preparedMatrix_le_one}
    \lean{TNLean.PEPS.PairEffect.Word.norm_preparedMatrix_tmul_le_one}
    \leanok
    \uses{def:peps_prepared_matrix}
    If $T$ is an allowed composition and $\|\eta_e\|=1$ in every
    slot, then $\|K_T(\eta)\|\le1$ in the Euclidean operator norm.
    In particular, if every $u_{e,r}$ and $v_{e,s}$ has norm one,
    then $\|K_T(a)\|\le1$ for every coordinate assignment $a$.
    The unit-vector families need not span their halfspaces.
    These are the coefficient contractions in
    \cite[proof of Theorem~5.2]{OpenAI2026PolynomialPEPS},
    \texttt{04-compression.tex}, lines 409--427.
\end{theorem}

\begin{proof}\leanok
    \uses{def:peps_source_inventory,def:peps_party_layout}
    The normalization of each source makes its preparation allowed,
    hence contractive. Composing with $T$ gives
    $\|TQ_{\eta,\mathcal H_{\rm in}}\|\le1$.
    Orthonormal coordinate maps are isometries, so the matrix in
    (\ref{eq:peps_prepared_matrix}) has the same operator norm.
    For the last assertion, use
    $\|u_{e,r}\otimes v_{e,s}\|=\|u_{e,r}\|\|v_{e,s}\|=1$.
\end{proof}

\begin{theorem}[Density coefficients constructed from a gate]
    \label{thm:peps_actual_gate_density}
    \lean{TNLean.PEPS.PairEffect.Word.exists_finite_density_expansion}
    \leanok
    \uses{def:peps_prepared_density_coefficient,def:peps_party_layout}
    Let $P$ and $\Xi$ be finite. For every $\xi\in\Xi$, let $F_\xi$
    be an allowed source-only composition on the parties $P$, with
    fixed input and output memories. Fix finite orthonormal bases of
    these memories and complex coefficients $c_\xi$, and let $M_\xi$
    be the matrix of $F_\xi$ in these bases.

    There is one ordered slot for each unordered pair of distinct
    parties, with common finite halfspaces $U_e=\C^{d_e}$ and
    $V_e=\C^{d'_e}$, normalized vectors
    $\eta_{\xi,e}\in U_e\otimes V_e$, and allowed compositions
    $T_\xi$ containing no sources such that
    $F_\xi=T_\xi Q_{\eta_\xi,\mathcal H_{\rm in}}$.
    Use the standard orthonormal bases of these halfspaces to define
    $K_\xi(a)$ as in
    Definition~\ref{def:peps_prepared_density_coefficient}, and put
    $\kappa_\xi(a)=\prod_e\eta_{\xi,e}(a_e)$.
    Each $K_\xi(a)$ has Euclidean operator norm at most one.

    These choices are independent of the input matrix $\rho$.
    For every complex input matrix $\rho$, writing
    $G=\sum_\xi c_\xi M_\xi$, one has
    \begin{align}
        G\rho G^\dagger
         & =\sum_{\xi,\zeta}c_\xi\overline{c_\zeta}
        \sum_{a,b}\kappa_\xi(a)\overline{\kappa_\zeta(b)}
        K_\xi(a)\rho K_\zeta(b)^\dagger.
        \label{eq:peps_actual_gate_density}
    \end{align}
    The halfspaces, source vectors, and coefficient matrices are
    obtained from the original compositions. This is the exact
    source-coordinate density identity preceding the random
    replacement in \cite[proof of Theorem~5.2]{OpenAI2026PolynomialPEPS},
    \texttt{04-compression.tex}, lines 233--355.
\end{theorem}

\begin{proof}\leanok
    \uses{thm:peps_finite_source_gate,thm:peps_prepared_gate_density,
        thm:peps_prepared_matrix_contraction}
    Theorem~\ref{thm:peps_finite_source_gate} constructs the common
    finite source spaces, their normalized vectors, and the exact
    source-free compositions $T_\xi$. Replacing every source half by
    a standard basis vector gives normalized elementary preparations;
    Theorem~\ref{thm:peps_prepared_matrix_contraction} therefore proves
    $\|K_\xi(a)\|\le1$.
    Taking matrices of $F_\xi=T_\xi Q_{\eta_\xi,\mathcal H_{\rm in}}$
    identifies $M_\xi$ with the prepared branch matrix.
    Apply Theorem~\ref{thm:peps_prepared_gate_density} to obtain
    (\ref{eq:peps_actual_gate_density}) for every $\rho$.
\end{proof}
